% Work/MAT — návrh revize hodiny 12, Co je množina.
% Rozpis: 1_rocnik/tematicke_plany/Rozpis_vyuky_MAT_1_rocnik.tex, hodina 12.
% Učebnice: ucebnice/01_Zakladni_poznatky/Mnoziny/Množiny.tex
%  (Množiny; Způsoby zápisu množin; Druhy množin; Rovnost množin).
% Realisticky.cz: metodická návaznost na lekci 010301 Množiny I;
%  lokální PDF není v tomto prostředí dostupné k úplnému porovnání.
% Záměrně použita vlastní formulace výkladu a samostatné úlohy;
% podrobné podmnožiny, operace a diagramy náleží následujícím hodinám.
\documentclass[11pt,a4paper]{article}
% Spolecny styl pro ucitelske pripravy MAT (Work/MAT)
% Verze 1.1, 2026-10-23
\NeedsTeXFormat{LaTeX2e}
\ProvidesPackage{priprava}[2026/10/23 v1.1 Sablona ucitelskych priprav MAT]

% Zakladni typografie a cestina
\RequirePackage[T1]{fontenc}
\RequirePackage[utf8]{inputenc}
\RequirePackage[czech]{babel}
\RequirePackage{lmodern}
\RequirePackage{microtype}
\RequirePackage{amsmath,amssymb}
\RequirePackage{xcolor}
\RequirePackage{enumitem}
\RequirePackage[most]{tcolorbox}
\RequirePackage{tabularx,booktabs,array}
\RequirePackage{fancyhdr}
\RequirePackage{needspace}
\RequirePackage{geometry}
\RequirePackage[hidelinks]{hyperref}

\geometry{a4paper,top=17mm,bottom=18mm,left=23mm,right=23mm}
\definecolor{notegray}{gray}{0.39}
\setlength{\parindent}{0pt}
\setlength{\parskip}{2.5pt plus 1pt minus 0.5pt}
\setlist[itemize]{topsep=2pt,itemsep=1pt,parsep=0pt,leftmargin=1.45em}
\setlist[enumerate]{topsep=2pt,itemsep=2pt,parsep=0pt,leftmargin=1.65em}
\setcounter{secnumdepth}{0}
\renewcommand{\arraystretch}{1.12}

% Zahlavi: vlevo rocnik | globalni tema, vpravo cislo: nazev hodiny
% Zapati zustava zcela prazdne.
\pagestyle{fancy}
\fancyhf{}
\renewcommand{\headrulewidth}{0.35pt}
\renewcommand{\footrulewidth}{0pt}
\setlength{\headheight}{14pt}
\newcommand{\nastavPripravu}[4]{%
  \fancyhead[L]{\small #1\enspace\textbar\enspace #2}%
  \fancyhead[R]{\small #3: #4}%
}

% Minimalisticke cernobile boxy. Kazda definice zvlast.
\tcbset{commonlesson/.style={enhanced,breakable=false,colback=white,colframe=black,
  colbacktitle=white,coltitle=black,boxrule=0.45pt,
  arc=0.8mm,outer arc=0.8mm,left=3mm,right=3mm,top=1.6mm,bottom=1.6mm,
  toptitle=1mm,bottomtitle=0.7mm,fonttitle=\bfseries,
  attach boxed title to top left={xshift=2.5mm,yshift=-1.15mm},
  boxed title style={colback=white,colframe=white,boxrule=0pt,left=1mm,right=1mm}}}
\newtcolorbox{zadaniBox}[1][]{commonlesson,title={Zadání},#1}
\newtcolorbox{prikladBox}[1][]{commonlesson,title={Příklad},#1}
\newtcolorbox{definiceBox}[1][]{commonlesson,title={Definice},#1}

% Spolecny box pro zadani a jeho reseni, bez mezery mezi castmi.
% Pouziti:
%   \begin{prikladReseniBox}
%     Zadani ulohy.
%     \tcblower
%     \begin{align*} ... \end{align*}
%   \end{prikladReseniBox}
% Spodni cast ma jemny sedy podklad a delici linku; lze v ni sazet
% vice radku matematiky. Obe casti maji jen jeden spolecny ramecek.
\tcbset{solutionlesson/.style={
  commonlesson,
  enhanced jigsaw,
  breakable,
  colbacklower=black!3,
  segmentation style={solid,draw=black!25,line width=0.35pt},
  fontlower=\small\color{notegray},
  before lower={\noindent\textbf{Řešení:}\enspace},
  leftlower=3mm,rightlower=3mm
}}
\newtcolorbox{prikladReseniBox}[1][]{solutionlesson,title={Příklad},#1}
\newtcolorbox{zadaniReseniBox}[1][]{solutionlesson,title={Zadání},#1}

% Pomucky pro autora pripravy: sedive poznamky; \reseni pro starsi pripravy
\newcommand{\poznamka}[1]{\par\begingroup\small\color{notegray}\itshape
  \textnormal{\textbf{Poznámka pro učitele.}}\enspace #1\par\endgroup}
\newcommand{\reseni}[1]{\par\begingroup\small\color{notegray}
  \textbf{Očekávané řešení:}\enspace #1\par\endgroup}
\newcommand{\blok}[1]{\needspace{5\baselineskip}\subsection*{#1}}
\newcommand{\vyrok}[1]{\uv{#1}}

\nastavPripravu{1. ročník}{Množiny}{12}{Co je množina}
\begin{document}
\blok{Motivace: co patří do souboru?}
Vybíráme čísla podle pravidla: \uv{všechna sudá přirozená čísla menší než 10}.
Dokážeme pro každé přirozené číslo rozhodnout, zda vyhovuje. Naproti tomu
\uv{pěkná čísla} nedává bez dalšího upřesnění jednoznačné pravidlo.

\begin{definiceBox}
\textbf{Množina} je v našem intuitivním pojetí soubor jednoznačně určených
objektů. O každém uvažovaném objektu umíme rozhodnout, zda do množiny patří.
\end{definiceBox}
\begin{definiceBox}
\textbf{Prvek množiny} je objekt, který do dané množiny patří.
Množiny obvykle označujeme velkými písmeny ($A,B,M$), prvky malými
písmeny ($a,b,x$) nebo konkrétními hodnotami.
\end{definiceBox}

Značení $a\in A$ znamená \uv{$a$ patří do $A$}, $a\notin A$ znamená
\uv{$a$ nepatří do $A$}. V množině $A=\{2,4,6,8\}$ platí
$4\in A$, $5\notin A$. \textbf{Na pořadí prvků nezáleží a opakování
prvku množinu nemění.}
\poznamka{Množina je prvotní pojem. Pro 1. ročník postačuje intuitivní
výklad; netvrdíme, že jde o úplnou axiomatickou definici. Před zadáním
formálního zápisu nech studenty samostatně vyslovit pravidlo výběru.}

\blok{Jak lze množinu zadat?}
\begin{definiceBox}
\textbf{Zadání výčtem}: uvedeme všechny prvky konečné množiny ve
složených závorkách, například $A=\{2,4,6,8\}$.
\end{definiceBox}
\begin{definiceBox}
\textbf{Zadání charakteristickou vlastností}: uvedeme obor, ze kterého
vybíráme, a podmínku, kterou musí prvky splňovat. Například
\[A=\{x\in\mathbb N\mid x\text{ je sudé a }x<10\}.\]
Svislou čáru čteme \uv{pro které platí}.
\end{definiceBox}

\poznamka{V projektu platí $\mathbb N=\{1,2,3,\ldots\}$;
$\mathbb Z$ obsahuje i nulu a záporná celá čísla.
U nekonečné množiny není zápis s \uv{$\ldots$} úplným konečným výčtem.
Rozlišuj zápis množiny a prostý seznam několika jejích příkladů.}

\begin{prikladReseniBox}[title={Příklad 1: vlastnost $\longrightarrow$ výčet}]
Zapiš výčtem $B=\{x\in\mathbb Z\mid -2\le x<3\}$.
Která z čísel $-2,\,2,\,3,\,\tfrac12$ do $B$ patří?
\tcblower
$B=\{-2,-1,0,1,2\}$. Patří $-2$ a $2$; $3$ je vyloučeno
ostrou mezí a $\tfrac12$ není celé číslo.
\end{prikladReseniBox}

\begin{prikladReseniBox}[title={Příklad 2: výčet $\longrightarrow$ vlastnost}]
Zapiš $C=\{-3,-1,1,3\}$ charakteristickou vlastností.
\tcblower
Například $C=\{x\in\mathbb Z\mid x\text{ je liché a }-3\le x\le3\}$.
Nestačí napsat jen $-3\le x\le3$: zahrnuli bychom i další čísla.
\end{prikladReseniBox}

\blok{Prázdná, jednoprvková, konečná a nekonečná množina}
\begin{definiceBox}
\textbf{Prázdná množina} neobsahuje žádný prvek. Značíme ji
$\varnothing$ nebo $\{\}$.
\end{definiceBox}
\begin{definiceBox}
\textbf{Jednoprvková množina} má právě jeden prvek. Například
$\{0\}$ obsahuje číslo nula; není prázdná.
\end{definiceBox}
\begin{definiceBox}
\textbf{Konečná množina} má konečně mnoho prvků; \textbf{nekonečná množina}
není konečná. Příklady: $\{1,3,5\}$ je konečná, $\mathbb N$ nekonečná.
\end{definiceBox}

\begin{prikladReseniBox}[title={Příklad 3: rozhodni a zdůvodni}]
Urči množiny $D=\{x\in\mathbb N\mid x<1\}$ a
$E=\{x\in\mathbb Z\mid -1<x<1\}$. Platí $D=E$?
\tcblower
$D=\varnothing$: žádné přirozené číslo není menší než $1$.
$E=\{0\}$: jediným celým číslem mezi $-1$ a $1$ je nula.
Proto $D\ne E$. Zápisy $0$, $\{0\}$ a $\varnothing$ nejsou zaměnitelné.
\end{prikladReseniBox}

\blok{Rovnost množin jako první kontrolní pravidlo}
\begin{definiceBox}
\textbf{Rovnost množin:} dvě množiny jsou si rovny právě tehdy,
když mají přesně stejné prvky. Pořadí ani opakování nerozhodují.
\end{definiceBox}
\begin{prikladReseniBox}[title={Příklad 4: co změní zápis?}]
Rozhodni, zda $\{1,2,3\}=\{3,2,1,2\}$ a zda
$\{0\}=\varnothing$.
\tcblower
První rovnost platí: v obou množinách jsou právě prvky $1,2,3$.
Druhá neplatí, protože levá množina má prvek $0$, pravá nemá žádný.
\end{prikladReseniBox}

\blok{Samostatné procvičení — základ hodiny}
\begin{zadaniReseniBox}[title={Procvičení S1}]
\begin{enumerate}[label=\alph*)]
\item Pro $M=\{-3,0,4\}$ doplň $\in$ nebo $\notin$:
$-3\ \ldots\ M$, $3\ \ldots\ M$, $0\ \ldots\ M$.
\item Vypiš prvky $P=\{x\in\mathbb Z\mid -3<x\le2\}$.
\item Zapiš $Q=\{2,4,6,8\}$ charakteristickou vlastností.
\item Rozhodni o rovnosti $\{0,1,1,2\}$ a $\{2,0,1\}$.
\item Která z množin $\{0\}$, $\varnothing$, $\mathbb Z$ je prázdná,
která konečná a která nekonečná?
\end{enumerate}
\tcblower
\begin{enumerate}[label=\alph*)]
\item $-3\in M$, $3\notin M$, $0\in M$.
\item $P=\{-2,-1,0,1,2\}$.
\item Například $Q=\{x\in\mathbb N\mid x\text{ je sudé a }x<10\}$.
\item Ano. Opakované číslo $1$ není dalším prvkem.
\item $\{0\}$ je konečná neprázdná, $\varnothing$ konečná prázdná,
$\mathbb Z$ nekonečná.
\end{enumerate}
\end{zadaniReseniBox}
\poznamka{Povinné jádro: motivace, význam prvku, dva zápisy, prázdná množina,
P1 a P3, S1 a)–d). P2, P4 a S1 e) dle tempa. Cílem je vědomě pracovat
s významem oboru, nikoli vyjmenovat co nejvíce definic.}

\blok{Volitelná rezerva}
\begin{zadaniReseniBox}[title={R1: stejná podmínka, různé obory}]
Urči výčtem, lze-li:
\begin{enumerate}[label=\alph*)]
\item $A=\{x\in\mathbb N\mid -1\le x<4\}$;
\item $B=\{x\in\mathbb Z\mid -1\le x<4\}$;
\item $C=\{x\in\mathbb R\mid -1\le x<4\}$.
\end{enumerate}
Které množiny jsou konečné?
\tcblower
$A=\{1,2,3\}$, $B=\{-1,0,1,2,3\}$,
$C=\langle-1;4)$ je interval reálných čísel.
Konečné jsou pouze $A,B$; $C$ je nekonečná.
\end{zadaniReseniBox}

\blok{Domácí procvičení}
\begin{zadaniReseniBox}[title={D1: procvič si zápisy}]
\begin{enumerate}[label=\alph*)]
\item Vypiš $A=\{x\in\mathbb Z\mid -4\le x<-1\}$.
\item Zapiš $B=\{1,3,5,7\}$ charakteristickou vlastností.
\item Urči $C=\{x\in\mathbb N\mid x<1\}$ a $D=\{x\in\mathbb Z\mid -1<x<1\}$.
\item Je množina všech záporných celých čísel konečná? Vysvětli.
\end{enumerate}
\tcblower
\begin{enumerate}[label=\alph*)]
\item $A=\{-4,-3,-2\}$.
\item Např. $B=\{x\in\mathbb N\mid x\text{ je liché a }x<8\}$.
\item $C=\varnothing$, $D=\{0\}$.
\item Ne, je nekonečná: pro každé záporné celé číslo lze nalézt další menší.
\end{enumerate}
\end{zadaniReseniBox}
\poznamka{Navazující hodina může zavést podmnožinu a operace;
v této přípravě je nevykládáme. Před zadáním konkrétních převzatých úloh
z Realisticky.cz ověř místní lekci 010301 a případně odkaž na čísla úloh
bez přepisu autorsky chráněných zadání.}
\end{document}
